% ==================================================================
% id:       func-003
% topic:    פונקציות
% subtopic: תחום הגדרה ותכונות של פונקציות
% solution: full
% ==================================================================
\textbf{הגדירו את המושגים הבאים:}
\begin{enumerate}\bMs{2}
\item
תחום ההגדרה הטבעי של נוסחה.
\begin{solbox}
תחום ההגדרה הטבעי של נוסחה
$f(x)$
הינו הקבוצה הגדולה ביותר של מספרים ממשיים שעליהם הנוסחה
$f(x)$
מוגדרת היטב.
\end{solbox}
\item
תמונה של פונקציה.
\begin{solbox}
תמונה של פונקציה
$f: A \rightarrow B$
היא אוסף האיברים של $B$ שיש להם מקור תחת הפונקציה. כלומר, אוסף האיברים
$y \in B$
כך שקיים
$x \in A$
עבורו
$f(x) =y$.
\end{solbox}
\item
פונקציה זוגית.
\begin{solbox}
פונקציה זוגית היא פונקציה
$f :D \rightarrow \rr$,
כאשר
$D \subseteq \rr$,
שבה לכל
$x \in D$
גם
$-x \in D$,
ומתקיים
$f(-x) = f(x)$.
\end{solbox}
\item
פונקציה אי-זוגית.
\begin{solbox}
פונקציה אי-זוגית  היא פונקציה
$f :D \rightarrow \rr$,
כאשר
$D \subseteq \rr$,
שבה לכל
$x \in D$
גם
$-x \in D$,
ומתקיים
$f(-x) =- f(x)$.
\end{solbox}
\eMs
\end{enumerate}
